% Appendices. \appendix is issued in main.tex.

\section{Configuration and Hyperparameters}
\label{app:config}

All experiments are driven by a single configuration file
(\texttt{configs/default.yaml}), the single source of truth referenced
throughout Sections~\ref{sec:architecture}--\ref{sec:training}. The tables
below reproduce the values used for the final run reported in this document.
The reward terms are not repeated here; they are given in
Table~\ref{tab:reward}.

\begin{table}[h]
\centering
\small
\begin{tabular}{lll}
\toprule
\textbf{Parameter} & \textbf{Value} & \textbf{Notes} \\
\midrule
Grid size            & $\gridsize\times\gridsize$ & fixed; never randomized (Invariant~1) \\
Action space         & $4$ (up/down/left/right)   & no ``stay'' action; masked at grid/obstacle edges \\
Vision metric        & Chebyshev                  & $(2\rvis{+}1)^2$ square window \\
Comm.\ metric        & Manhattan                  & range-limited pairwise map fusion \\
Global channels      & $6$                        & Visited, Obstacle, Trajectory, Target, Own\_Pos., Broken \\
Local channels       & $5$                        & cropped $13\times13$ egocentric patch \\
Context dimension    & $5$                        & $[\rvis, \rcomm, \nagents, \text{agent\_id}, \hat n]$ \\
Discount $\gamma$    & $0.97$                     & effective horizon for search \\
\bottomrule
\end{tabular}
\caption{Environment and observation constants.}
\label{tab:app-env}
\end{table}

\begin{table}[h]
\centering
\small
\begin{tabular}{lll}
\toprule
\textbf{Factor} & \textbf{Range (final run)} & \textbf{Normalization denom.} \\
\midrule
Vision radius $\rvis$      & $[1, 6]$ (integer)     & $6$ \\
Comm.\ range $\rcomm$      & $[2, 12]$ (integer)    & $12$ \\
Team size $\nagents$       & $[1, 4]$ (integer)     & $4$ (also used for agent\_id, $\hat n$) \\
Obstacle density $\rho$    & $[0.10, 0.30]$ (float) & environment-level (not a context entry) \\
Fault prob.\ $\pfault$     & $[0.0, 0.003]$ (float) & environment-level (not a context entry) \\
\bottomrule
\end{tabular}
\caption{Per-episode domain-randomization ranges, sampled uniformly at every
reset (Section~\ref{subsec:train-dr}). Context factors are normalized by their
maxima (Section~\ref{subsec:arch-context}); density and fault probability are
deliberately \emph{not} exposed to the network.}
\label{tab:app-dr}
\end{table}

\begin{table}[h]
\centering
\small
\begin{tabular}{lll}
\toprule
\textbf{Phase} & \textbf{Episodes} & \textbf{Overrides} \\
\midrule
\texttt{test}               & $0$ (skipped)  & staging placeholder; narrower \texttt{dr\_override} \\
\texttt{mas\_50k\_dr\_faults} & $50{,}000$   & \texttt{max\_steps}${=}200$, \texttt{n\_agents}${=}4$ (buffer sizing) \\
\bottomrule
\end{tabular}
\caption{Curriculum. The final policy is trained from scratch in the single
main phase; the zero-episode debug phase is kept only for staging
(Section~\ref{subsec:train-curriculum}). The phase's \texttt{n\_agents} sets
the initial allocation and replay-buffer width; team size is still resampled
in $[1,4]$ every episode from the ranges of Table~\ref{tab:app-dr}.}
\label{tab:app-curriculum}
\end{table}

\begin{table}[h]
\centering
\small
\begin{tabular}{lll}
\toprule
\textbf{Hyperparameter} & \textbf{Value} & \textbf{Notes} \\
\midrule
Optimizer                 & Adam           & weight decay $10^{-6}$ \\
Learning rate             & $10^{-4}\!\to\!5\times10^{-4}\!\to\!10^{-6}$ & linear warmup ($10\%$) then cosine; per episode \\
Gradient clip             & $1.0$          & global norm \\
Discount $\gamma$         & $0.97$         & --- \\
$n$-step return           & $3$            & per-agent windows; team credit needs $n\!\ge\!2$ \\
Loss                      & Huber (smooth $L_1$) & Double-DQN target \\
Replay buffer             & $2\times10^5$  & uniform; not serialized in checkpoints \\
Batch size                & $16$           & --- \\
Update cadence            & every $4$ agent-turns & one gradient step \\
Target sync               & every $500$ gradient steps & hard copy \\
$\varepsilon$ schedule    & $1.0\!\to\!0.05$, $\times0.9995$/episode & floor at $\approx$ ep.\ $6{,}000$ \\
FC head widths            & $1536\!\to\!768\!\to\!384$ & dueling head $Q=V+A-\overline{A}$ \\
\bottomrule
\end{tabular}
\caption{Agent and training hyperparameters (Sections~\ref{subsec:arch-network},
\ref{subsec:arch-learning}, \ref{subsec:train-schedules}). Network
normalization is LayerNorm/GRN only, for batch-size-independent inference.}
\label{tab:app-agent}
\end{table}

\begin{table}[h]
\centering
\small
\begin{tabular}{lll}
\toprule
\textbf{Setting} & \textbf{Value} & \textbf{Purpose} \\
\midrule
In-training eval interval  & every $1{,}000$ episodes & checkpoint selection \\
In-training eval instances & $20$ fixed seeds         & success-dominant rule \\
In-training eval $\varepsilon$ & $0.05$               & near-greedy \\
Periodic checkpoint save   & every $2{,}500$ episodes & learning-curve axis \\
Parallel watcher eval      & $200$ seeded episodes/checkpoint & crash-safe second opinion \\
OFAT sweep seeds           & $500$ per point          & Section~\ref{sec:experiments} \\
Learning-curve / heatmap seeds & $300$ per point      & Section~\ref{sec:experiments} \\
\bottomrule
\end{tabular}
\caption{Evaluation protocol (Sections~\ref{subsec:train-infra},
\ref{subsec:exp-setup}).}
\label{tab:app-eval}
\end{table}

\FloatBarrier

\section{Repository Guide}
\label{app:repo}

The code is organized as a small set of packages behind three entry points.

\paragraph{Module map.}
\begin{itemize}[leftmargin=1.5em]
  \item \texttt{env/} --- the environment package: \texttt{grid\_env.py}
    (the \texttt{DroneSearchEnv}: generation with BFS reachability, sequential
    \texttt{step\_agent}, reward, beacon detection, communication fusion) and
    \texttt{utils.py} (BFS helpers and the eval-only \texttt{auto\_nav\_action}).
  \item \texttt{agents/} --- \texttt{networks.py} (GlobalMapCNN, LocalCNN,
    dueling head) and \texttt{dqn\_agent.py} (Double-DQN agent: action
    selection, learning step, checkpoint I/O).
  \item \texttt{training/} --- \texttt{train.py} (curriculum loop, domain
    randomization, $n$-step windows, in-training evaluation).
  \item \texttt{gui/} --- the renderer for the \texttt{play}, \texttt{eval},
    and \texttt{simulate} modes (fault slider, auto-nav toggle).
  \item \texttt{testing/} --- \texttt{evaluate\_policy.py} (the OFAT / fault /
    heatmap / epoch harness) and \texttt{analysis.ipynb} (figure generation).
  \item Top level --- \texttt{main.py} (all runtime modes),
    \texttt{eval\_checkpoints.py} (batch and \texttt{--watch} checkpoint
    ranking), and \texttt{configs/default.yaml} (single source of truth).
\end{itemize}

\paragraph{Environment setup.}
\begin{verbatim}
conda env create -f environment.yaml   # env "rl-drone": python 3.11, CUDA 11.8
conda activate rl-drone
# or, without conda:  pip install -r requirements.txt
\end{verbatim}

\paragraph{Reproducing the results.}
\begin{verbatim}
# Training (curriculum from configs/default.yaml; --resume to continue)
python main.py --mode train

# Qualitative inspection
python main.py --mode eval --checkpoint checkpoints/best.pt \
                                          [--fault-prob P] [--auto-nav]
python main.py --mode simulate            # GUI: fault slider, auto-nav toggle

# The experiments of Section 7 (writes testing_results/{csv,plots})
python testing/evaluate_policy.py --axis all      --seeds 500
python testing/evaluate_policy.py --axis fault    --seeds 500 --auto-nav
python testing/evaluate_policy.py --axis heatmaps --seeds 300
python testing/evaluate_policy.py --axis epoch    --seeds 300
# then run testing/analysis.ipynb to regenerate the figures
\end{verbatim}

\FloatBarrier

\section{Additional Figures}
\label{app:figures}

For completeness this appendix collects four figures referenced from the body but
not shown there. Figures~\ref{fig:app-vision} and~\ref{fig:app-maxsteps} give the
two one-factor sweeps summarized in Table~\ref{tab:generalization} --- the
vision-radius and step-budget axes --- with the six panels of each following the
layout of Figures~\ref{fig:sweep-nagents} and~\ref{fig:sweep-density}.
Figures~\ref{fig:schedules}--\ref{fig:epoch} then document the training run of
Section~\ref{sec:training}: its two per-episode schedules, the in-training
evaluation that drives checkpoint selection, and the checkpoint-level learning
curve discussed in Section~\ref{subsec:exp-generalization}.

\begin{figure}[h]
\centering
\includegraphics[width=\textwidth]{fig_vision.png}
\caption{Sensitivity to vision radius $\rvis$. Coverage rises monotonically
with sight (top right), but success saturates around $\rvis = 3\text{--}4$:
past a point the bottleneck is deciding where to go, not seeing farther. The
OOD radii $\rvis = 7, 8$ (shaded) degrade only slightly.}
\label{fig:app-vision}
\end{figure}

\begin{figure}[h]
\centering
\includegraphics[width=\textwidth]{fig_max_steps.png}
\caption{Sensitivity to the step budget $T$. Success rises with diminishing
returns as the budget grows; a larger budget only rescues episodes that were
already close to succeeding. No OOD region is defined for this axis (the
budget is not a randomized factor).}
\label{fig:app-maxsteps}
\end{figure}

\begin{figure}[t]
\centering
\includegraphics[width=0.92\textwidth]{fig_schedules.png}
\caption{The two per-episode schedules of the final $50{,}000$-episode run.
\emph{(a)} $\varepsilon$ decays multiplicatively to its $0.05$ floor at episode
$5{,}990$; \emph{(b)} the learning rate warms up over the first $10\%$ of
episodes to $5\times10^{-4}$, then anneals cosine-wise to $10^{-6}$. The two are
deliberately out of phase: exploration has essentially finished annealing while
the learning rate is still near its peak, so the bulk of the high-learning-rate
budget is spent on a near-greedy --- and therefore task-relevant --- state
distribution.}
\label{fig:schedules}
\end{figure}

\begin{figure}[t]
\centering
\includegraphics[width=0.86\textwidth]{fig_training_eval.png}
\caption{In-training greedy evaluation of the same run, executed every $1{,}000$
episodes on the fixed bank of $20$ seeded instances that drives checkpoint
selection; \emph{(a)} success rate and \emph{(b)} mean return, best values
circled. These levels are not comparable with those of
Section~\ref{sec:experiments}: this bank samples the \emph{full} randomization
ranges --- solo drones, $\rvis{=}1$, faults active --- whereas
Section~\ref{sec:experiments} measures at the reference operating point. The
$5$-point sawtooth in \emph{(a)} is the quantization of a $20$-instance bank,
which is the noise the strict-improvement selection rule is designed not to
chase. Both panels flatten by around episode $7{,}000$.}
\label{fig:training-eval}
\end{figure}

\begin{figure}[t]
\centering
\includegraphics[width=0.86\textwidth]{fig_epoch.png}
\caption{Learning curve: every periodic checkpoint of the final run,
evaluated greedily at the reference operating point ($300$ seeds). Success
plateaus near $85\%$ while SPL keeps improving, indicating late refinement of
path quality rather than of find rate.}
\label{fig:epoch}
\end{figure}
